\begin{table}[htbp]
\centering
\caption{Accounting components for selected focal coalition vectors}
\label{tab:accounting-focal-cases}
\scriptsize
\setlength{\tabcolsep}{3.0pt}
\begin{tabularx}{\textwidth}{@{}lXrrrrrL{1.4cm}@{}}
\toprule
Domain & Election/case & Vote \% & Seats & \(d_C\) & \(A_C\) & \(B_C\) & Pattern \\
\midrule
Cabinet & Cabinet 2014/2016.2 & 46.54 & 259 & 20.271 & 15.742 & 4.529 & A+, B+ \\
Cabinet & Cabinet 2014/2017.1 & 49.91 & 266 & 9.952 & 11.268 & -1.316 & A+, B- \\
Cabinet & Cabinet 2018/2021.3/2022.1 & 47.25 & 257 & 14.624 & 19.069 & -4.445 & A+, B- \\
Cabinet & Cabinet 2022/2023.1 & 48.89 & 263 & 12.204 & 13.196 & -0.992 & A+, B- \\
Ideological & Ideological 2014/PTB-PR & 47.59 & 257 & 12.857 & 11.621 & 1.236 & A+, B+ \\
Ideological & Ideological 2022/MDB-UNIÃO & 49.98 & 265 & 8.591 & -1.292 & 9.883 & A-, B+ \\
Ideological & Ideological 2022/PP-PL & 45.35 & 258 & 25.355 & 22.821 & 2.534 & A+, B+ \\
\bottomrule
\end{tabularx}
\begin{minipage}{0.96\linewidth}
\footnotesize Notes: The two 2018 cabinet periods share one election-space numerical vector and therefore appear once. The ideological rows are focal endpoint-minimal cases; all minimal connected ideological inversions are reported in Table~\ref{tab:minimal-intervals}. The identities are descriptive, not causal.
\end{minipage}
\end{table}
